\chapter{Grothendieck groups and composition factors}
\label{chap:library:kzero}

The Grothendieck group assigns a class $[V]$ to each representation, with
$[V]=[U]+[W]$ for every short exact sequence $0\to U\to V\to W\to0$.
For finite dimensional representations, the Jordan--H\"older theorem shows
that the classes of simple modules form a basis. Thus composition
multiplicities determine the class of a representation.

In the parts of this chapter concerning representations, $G$ is a finite
group and $k$ is a field. No splitting or characteristic assumption is needed
until the final section on Brauer characters. Write $kG$ for the group algebra,
$\operatorname{FDRep}_k(G)$ for the category of finite dimensional
$k$-representations, and $K_0(kG)$ for its Grothendieck group defined by short exact sequences.

\section{The Grothendieck group and its universal property}
\label{lib:kzero:presentation}

\begin{definition}
Let $\mathcal C$ be an abelian category whose isomorphism classes of
objects form a set.
Let $F(\mathcal C)$ be the free abelian group on the isomorphism classes
of its objects.
The Grothendieck group defined by short exact sequences is
\[
 K_0(\mathcal C)=F(\mathcal C)/R(\mathcal C),
\]
where
\[
 R(\mathcal C)=
 \left\langle e_V-e_U-e_W:
 0\longrightarrow U\longrightarrow V\longrightarrow W\longrightarrow0
 \text{ is exact}\right\rangle.
\]
The symbol $e_X$ denotes the generator corresponding to the isomorphism
class of $X$. Its image in the quotient is denoted by $[X]$.
\end{definition}

For categories of modules, this definition is recalled in
\cite[Appendix, p.~163]{Serre1977}.
Isomorphic objects have the same class. Each defining relation gives
$[V]=[U]+[W]$. The short exact
sequence with all three terms zero implies
$[0]=[0]+[0]$, hence
$[0]=0$. The split sequence for $U\oplus W$ gives
$[U\oplus W]=[U]+[W]$. Exact relations also apply to extensions that do not
split, which is the feature needed for modular representations.

\begin{proposition}[Universal property]
\label{lib:kzero:universal}
Let $A$ be an abelian group. An isomorphism invariant assignment $X\mapsto a_X$
defines a homomorphism $K_0(\mathcal C)\to A$ with $[X]\mapsto a_X$ if and
only if
\[
 a_V=a_U+a_W
\]
for every short exact sequence $0\to U\to V\to W\to0$. Such a homomorphism is
unique. Consequently two homomorphisms from $K_0(\mathcal C)$ agree if they
agree on the classes $[X]$.
\end{proposition}
\begin{proof}
The assignment defines a unique homomorphism $F(\mathcal C)\to A$ by the
universal property of the free abelian group. It vanishes on
$e_V-e_U-e_W$ precisely when the stated identity holds. Vanishing on the
generators of $R(\mathcal C)$ is equivalent to vanishing on that subgroup,
so the homomorphism descends to the quotient. Conversely every homomorphism
on the quotient respects its defining relations. Object classes generate
the quotient, which proves uniqueness.
\end{proof}

\begin{implementation}
\module{ModularRep.ExactGrothendieckGroup} constructs this presentation
from Mathlib's skeleton, free abelian group and short exact complexes. Here
\lean{FreeClassGroup} denotes the free abelian group,
\lean{shortExactRelationSubgroup} its subgroup of relations, and
\lean{KZero} the quotient. The identities for isomorphic objects and
short exact sequences are \lean{classOf_iso} and
\lean{classOf_middle_eq}. The universal property is \lean{lift}, with
\lean{lift_classOf} giving its value on an object class and
\lean{hom_ext} expressing uniqueness. The elements of {\small\texttt{Skeleton C}}
are isomorphism classes of objects.
These definitions also apply to categories with zero morphisms, using
short exact complexes $U\to V\to W$. The statements above concern
abelian categories.
\end{implementation}

\begin{proposition}[Maps induced by exact functors]
\label{lib:kzero:functor}
Let $F:\mathcal C\to\mathcal D$ be an exact functor between abelian
categories, that is, an additive functor that preserves short exact
sequences. Then $F$ induces a homomorphism
\[
 K_0(F):K_0(\mathcal C)\longrightarrow K_0(\mathcal D),
 \qquad [X]\longmapsto[F(X)].
\]
\end{proposition}
\begin{proof}
The hypotheses ensure that the image of a short exact sequence is short
exact. Thus the free homomorphism induced by $F$ sends every defining
relation to a defining relation for $K_0(\mathcal D)$. It therefore induces a
homomorphism on the quotients.
\end{proof}

\begin{implementation}
The induced homomorphism is \lean{ExactGrothendieckGroup.map}.
Its hypotheses require preservation of zero morphisms, finite limits and
finite colimits. For abelian categories, these are equivalent to exactness.
The construction uses
\lean{freeMap_shortExactRelation} and
\lean{map_shortExactRelationSubgroup}. The formula on object classes is
\lean{map_classOf}. No choice of bases is involved.
\end{implementation}

\section{Composition factors and the Jordan--H\"older theorem}
\label{lib:kzero:finite-length}

Let $R$ be a ring and $M$ a left $R$-module of finite length. A composition series of $M$ is a chain
\[
 0=M_0<M_1<\cdots<M_r=M
\]
whose successive quotients $M_{i+1}/M_i$ are simple. We also consider
composition series between two fixed submodules $A\leq B$ of $M$, with
first term $A$ and last term $B$. For an account of composition series and
the Jordan--H\"older theorem, see, for example,
\cite[\S1.1, especially Theorem~1.1.4]{Benson1998}.

\begin{definition}
Write $\operatorname{Mod}R$ for the category of all left $R$-modules in
the chosen universe. Recall that $F(\operatorname{Mod}R)$ is the free
abelian group on their isomorphism classes.
For a composition series $s$ of $M$, its formal factor sum is
\[
 J(s)=\sum_{i=0}^{r-1} e_{M_{i+1}/M_i}
 \in F(\operatorname{Mod}R),
\]
and the coefficient $m_X(s)$ at an isomorphism class $X$ counts the factors
isomorphic to $X$.
\end{definition}

\begin{proposition}[Independence of composition multiplicities]
\label{lib:kzero:series-independent}
Any two composition series $s$ and $t$ between the same two submodules
of $M$ have the same simple factors up to isomorphism, counted with
multiplicity.
Thus $J(s)=J(t)$ and $m_X(s)=m_X(t)$ for every module
isomorphism class $X$. In particular, a composition series from $0$ to a
module of finite length defines $J(M)$ and $m_X(M)$ independently of
the choice of series.
\end{proposition}
\begin{proof}
The Jordan--H\"older theorem gives a bijection between the factor indices of
the two series and an isomorphism between corresponding factors. After
passing to isomorphism classes, reindexing the finite sum gives the result.
This proves equality in the free abelian group, before imposing exact relations.
\end{proof}

\begin{implementation}
\module{ModularRep.JordanHolderCoordinates} defines
\lean{compositionFactor}, \lean{compositionFactorIsoSum}, and
\lean{compositionCoefficient}. Each quotient is formed inside its
numerator: the denominator in
\lean{compositionFactor} is the inverse image of the preceding submodule
under the numerator's inclusion into $M$. Simplicity follows from the
covering relation between adjacent submodules in
\lean{compositionFactor_isSimple}.
The proof of \lean{compositionFactorIsoSum_eq} applies Mathlib's
\lean{CompositionSeries.jordan_holder} and the corresponding module
isomorphisms. The formula
\lean{compositionCoefficient_eq_sum_ite} counts the factors in each
isomorphism class.
\end{implementation}

Every finite dimensional representation has finite length as a module over
its monoid algebra, even when the monoid is infinite. Indeed it is both
Noetherian and Artinian as a vector space. A chain of submodules over the
monoid algebra is also a chain of vector subspaces, so both chain conditions
pass to that module structure. In
\module{ModularRep.FDRepFiniteLength},
\lean{asModule_isFiniteLength} applies Mathlib's criteria for finite length
and its results on restriction of scalars. The theorem
\lean{exists_compositionSeries} then applies Mathlib's equivalence between
finite length and the existence of a composition series from $0$ to the whole module.
For finite $G$, all simple $kG$-modules are finite dimensional, as shown
in Section~\ref{lib:kzero:simple-basis}.

\section{Splicing series in a short exact sequence}
\label{lib:kzero:splicing}

\begin{proposition}[Additivity of composition factors]
\label{lib:kzero:factor-additivity}
For a short exact sequence of finite length left $R$-modules
\[
 0\longrightarrow U\xrightarrow{f}V\xrightarrow{\pi}W\longrightarrow0,
\]
the formal factor sums satisfy $J(V)=J(U)+J(W)$. Consequently
$m_X(V)=m_X(U)+m_X(W)$ for every module isomorphism class $X$.
\end{proposition}
\begin{proof}
Choose composition series
$0=U_0<\cdots<U_a=U$ and $0=W_0<\cdots<W_b=W$.
The images $f(U_i)$ form a composition series from $0$ to $\im f$, since
$f$ is injective. The inverse images $\pi^{-1}(W_j)$ form one from
$\ker \pi$ to $V$, since $\pi$ is surjective. Exactness identifies the meeting
terms $\im f=\ker \pi$, so the two chains concatenate to
\[
 0=f(U_0)<\cdots<f(U_a)=\pi^{-1}(W_0)
 <\cdots<\pi^{-1}(W_b)=V.
\]
The factor isomorphisms are
\[
 U_{i+1}/U_i\simeq f(U_{i+1})/f(U_i),
 \qquad
 \pi^{-1}(W_{j+1})/\pi^{-1}(W_j)\simeq W_{j+1}/W_j.
\]
The first is induced by the injective map $f$. For the second,
restrict $\pi$ to $\pi^{-1}(W_{j+1})$. The induced map on quotients is
surjective and has zero kernel, since the inverse image of $W_j$ is
$\pi^{-1}(W_j)$. Thus the concatenated series has
factor sum $J(U)+J(W)$. Jordan--H\"older compares it with any chosen full
series of $V$.
\end{proof}

\begin{implementation}
The proof in \module{ModularRep.FDRepJordanHolderKZero} constructs these
chains as \lean{mapCompositionSeries} and \lean{comapCompositionSeries}.
The factor isomorphisms are
\lean{quotientMapEquivOfInjective} and
\lean{quotientComapEquivOfSurjective}. Concatenation uses Mathlib's
composition series operation \lean{smash}, and the local theorem
\lean{compositionFactorIsoSum_smash} expresses its factor sum as the sum
of the two factor sums. This gives
\lean{compositionFactorIsoSum_eq_add_of_exact}.

For representations, \lean{fdRepJordanHolderSum_shortExact} applies the
exact underlying $kG$-module functor and uses the injection, surjection,
and exactness of the resulting linear maps. The universal property gives
\lean{fdRepJordanHolderKZeroHom}, with values in the free abelian group
on module isomorphism classes.
\end{implementation}

\section{The basis of simple modules}
\label{lib:kzero:simple-basis}

Let $\mathcal S_k(G)$ denote the set of isomorphism classes of simple left
$kG$-modules. Every such module is cyclic: any nonzero vector generates a
nonzero submodule and hence the whole module. Since $G$ is finite, $kG$ is
finite dimensional over $k$, so every simple module belongs to
$\operatorname{FDRep}_k(G)$.

For an account of the following standard description of the Grothendieck
group, see \cite[\S14.1, Proposition~40]{Serre1977}.

\begin{theorem}[Basis of simple module classes]
\label{lib:kzero:basis-theorem}
The map
\[
 K_0(kG)\longrightarrow\bigoplus_{S\in\mathcal S_k(G)}\Z e_S,
 \qquad [V]\longmapsto\sum_S m_S(V)e_S
\]
is an isomorphism of abelian groups. Its inverse sends $e_S$ to $[S]$.
\end{theorem}
\begin{proof}
All factors in $J(V)$ are simple, so its sum lies in the subgroup freely
generated by $\mathcal S_k(G)$. The additivity just proved gives the
homomorphism in the statement. In the other direction, the free abelian
group admits a homomorphism $e_S\mapsto[S]$.

For a composition series from $0$ to $V$, each step has a short exact sequence
\[
 0\longrightarrow V_i\longrightarrow V_{i+1}
 \longrightarrow V_{i+1}/V_i\longrightarrow0.
\]
Starting from $[0]=0$ and successively applying the defining relations gives
\[
 [V]=\sum_{i=0}^{r-1}[V_{i+1}/V_i].
\]
Thus the composite back to $K_0(kG)$ fixes every object class and hence the
whole group. For a simple module $S$, the series $0<S$ has the single
factor $S$. The other composite therefore
fixes every generator $e_S$. The maps are inverse.
\end{proof}

\begin{implementation}
In \module{ModularRep.FDRepSimpleClassKZero},
\lean{SimpleModuleClass} is the subtype of the module skeleton represented
by simple modules. The homomorphism \lean{simpleClassProjection} fixes the
generators represented by simple modules and sends the others to zero.
Composing it with the homomorphism defined by composition multiplicities gives
\lean{fdRepSimpleJordanHolderHom}.

The reverse map is \lean{simpleClassToFDRepKZero}.
The proof first sums the step relations in the category
{\small\texttt{FGModuleCat k[G]}} of finitely generated $kG$-modules using
\lean{compositionSeriesFGKZero_decomposition}, and then applies
\lean{FDRepGroupAlgebraEquivalence.equivalenceFGModule}.
This is a restriction of Mathlib's equivalence between representations
and modules over the monoid algebra. Finiteness of $G$ ensures that finite
generation over $kG$ implies finite dimension over $k$. The resulting
isomorphism is
\lean{fdRepKZeroEquivSimpleModuleClassGroup}.
\end{implementation}

\begin{example}[Classes of extensions of simple modules]
\label{lib:kzero:extension-example}
Suppose $0\to S\to V\to T\to0$ is short exact, where $S$ and $T$ are
simple. Then $[V]=[S]+[T]$. If $S\not\simeq T$, the coefficients of
$[S]$ and $[T]$ are one and all other coefficients are zero. If $S\simeq T$,
the coefficient of $[S]$ is two and all other coefficients are zero.
Thus any two such extensions have the same class in $K_0(kG)$, even when
their middle terms are not isomorphic.
\end{example}

For virtual classes, coefficients may be negative. Equality in $K_0(kG)$
means equality of all these integer coefficients. For a representation
$V$, the coefficient of $[S]$ is $m_S(V)$, its composition multiplicity.
This is the identity
\lean{FDRepJordanHolderKZero.fdRepJordanHolderCoordinate_classOf}.

\section{Finiteness of the set of simple classes}
\label{lib:kzero:simple-finite}

\begin{proposition}[Finiteness of the simple module classes]
\label{lib:kzero:finite-simples-theorem}
If the left regular module of a ring $R$ has a full finite composition
series, every simple left $R$-module occurs among its factors. In particular
$\mathcal S_k(G)$ is finite for a finite group $G$ and any field $k$.
\end{proposition}
\begin{proof}
A simple module $S$ is isomorphic to $R/I$ for a maximal left ideal $I$:
the map $R\to S$ sending $r$ to $rs$ is onto for any nonzero $s\in S$.
Since $R$ has finite length, so does $I$. Append $R$ to a
composition series of $I$. The resulting series of $R$ has final
factor $R/I\simeq S$. The Jordan--H\"older theorem identifies this
factor with a factor of any prescribed composition series from $0$ to $R$.

The finite index set of the factors therefore maps onto all simple module
classes. For $R=kG$, finite dimension over $k$ makes the regular module
both Noetherian and Artinian, so it has a full finite composition series.
\end{proof}

\begin{implementation}
\module{ModularRep.SimpleModuleClassFinite} proves
\lean{simpleClass_appears_in_compositionSeries} using a series whose
last factor is the quotient by a maximal left ideal.
\lean{finite_simpleModuleClass_of_compositionSeries} applies finiteness
to the resulting surjection. The result
\lean{finiteSimpleModuleClassGroupAlgebra} uses the chain conditions
that follow from finite dimension of $kG$.
\end{implementation}

Proposition~\ref{lib:kzero:finite-simples-theorem}, together with
Theorem~\ref{lib:kzero:basis-theorem}, proves that $K_0(kG)$ is free
of finite rank over $\Z$, for every field $k$.

\section{The Grothendieck group of all modules}
\label{lib:kzero:swindle}

\begin{proposition}[Eilenberg swindle]
\label{lib:kzero:swindle-theorem}
For a ring $R$, the exact Grothendieck group of the category of all left
$R$-modules in the chosen universe is zero.
\end{proposition}
\begin{proof}
For any $R$-module $M$, let $T=\prod_{n\geq0}M$. The shift isomorphism
\[
 M\oplus T\longrightarrow T,
 \qquad (x,f)\longmapsto(x,f(0),f(1),\ldots)
\]
is $R$-linear. The split short exact sequence
$0\to M\to M\oplus T\to T\to0$ gives $[M\oplus T]=[M]+[T]$.
The isomorphism gives $[M\oplus T]=[T]$, so cancellation yields $[M]=0$.
Since object classes generate the group, every element is zero.
\end{proof}

\begin{implementation}
\module{ModularRep.ModuleCategoryKZeroTrivial} uses the countable product
{\small\texttt{Nat -> M}}. Its results are \lean{classOf_eq_zero} and
\lean{kZero_subsingleton}. The homomorphism of composition factors
\lean{fdRepJordanHolderKZeroHom} takes values in
{\small\texttt{FreeClassGroup (ModuleCat k[G])}}, the free abelian group on
isomorphism classes of all $kG$-modules. Imposing all short exact relations would give
{\small\texttt{KZero (ModuleCat k[G])}}, which is zero by this argument.
\end{implementation}

For a nonzero $kG$-module $M$, the product $\prod_{n\geq0}M$ is not finite
dimensional over $k$, so the argument does not apply in
$\operatorname{FDRep}_k(G)$.
This is why the restriction to finite dimensional representations is
necessary in Theorem~\ref{lib:kzero:basis-theorem}.

\section{The Brauer character homomorphism}
\label{lib:kzero:characters}

Now let $\ell$ be prime, let $k$ be algebraically closed of characteristic
$\ell$, and let $K$ be a field of characteristic zero. Fix a finite root
correspondence $\iota$ for $G$ as in Chapter~\ref{chap:library:brauer}.
Let $\operatorname{CF}_{\ell'}(G,K)$ be the additive group of $K$-valued
class functions on the $\ell$-regular elements of $G$.

\begin{proposition}[Injectivity of the Brauer character homomorphism]
\label{lib:kzero:brauer-injective}
Taking Brauer characters induces an injective homomorphism
\[
 \operatorname{br}_{\iota}:K_0(kG)\longrightarrow
 \operatorname{CF}_{\ell'}(G,K),
 \qquad [V]\longmapsto\varphi_V.
\]
\end{proposition}
\begin{proof}
Isomorphism invariance and short exact additivity of Brauer characters
allow the universal property to define the homomorphism. Under the basis
of Theorem~\ref{lib:kzero:basis-theorem}, it sends $e_S$ to the irreducible
Brauer character $\varphi_S$. The functions $\varphi_S$ are pairwise distinct by
Proposition~\ref{lib:brauer:separation}, and linearly independent over
$K$ by Theorem~\ref{lib:brauer:independence}. Linear independence forces
each coefficient in an integer relation to vanish in $K$. Since
$\operatorname{char}K=0$, the map $\Z\to K$ is injective, so the
coefficients vanish in $\Z$.
\end{proof}

\begin{implementation}
\module{ModularRep.BrauerCharacterKZero} constructs
\lean{FDRepSimpleClassKZero.brauerCharacterKZeroHom} using exact
additivity. The theorem \lean{brauerCharacterKZeroHom_injective}
uses the basis of simple module classes and
\lean{irreducibleBrauerCharacterLinearIndependence_of_rootEmbedding}.
The proof restricts scalars from $K$ to $\Z$ and uses the injectivity of
the corresponding map taking a finite integer linear combination of characters.
\end{implementation}

Thus representations with equal Brauer characters have the same
composition multiplicities. In Chapter~\ref{chap:library:reduction},
character compatibility then identifies the Grothendieck classes of
different lattice reductions and proves independence of the chosen lattice.
